% !TeX root = ../../../金融经济学_王江_讲义及习题.tex
% 来源：杨云雯第12章批注，定理12.7旁的权重空间草图。
\begin{tikzpicture}[>=Stealth,font=\small,
 frontier/.style={blue!65!black,very thick,<->}]
 \begin{scope}[yshift=.3cm]
  \filldraw[fill=blue!5,draw=gray!65]
   (0,0)--(3.95,0)--(5.65,2.95)--(1.7,2.95)--cycle;
  \node[above] at (3.63,2.98) {预算平面 $\iota^{\mathsf T}z=1$};
  \draw[frontier] (.7,.6)--(5,2.32);
  \fill (1.3,.84) circle(1.7pt) node[above left] {$z_0$};
  \fill (2.75,1.42) circle(1.7pt) node[above left] {$z_{\mathrm{MVP}}$};
  \fill (4.2,2) circle(1.7pt) node[above left] {$z_1$};
  \draw[->,blue!65!black] (2.7,.9)--(3.9,1.38) node[midway,below right] {$h$};
  \node at (2.83,-.65) {(a) 预算平面内的视图};
 \end{scope}

 \begin{scope}[xshift=9.55cm,yshift=1.14cm,
  x={(-1.65cm,-.61cm)},y={(2cm,-.42cm)},z={(0cm,2.45cm)}]
  % 投影方向和预算平面边界与图 12.S11(b) 相同。
  \fill[blue!6]
   (1.46,-.26,-.20)--(-.26,1.46,-.20)--(-.99,.73,1.26)--(.73,-.99,1.26)--cycle;
  \filldraw[fill=blue!16,draw=blue!40] (1,0,0)--(0,1,0)--(0,0,1)--cycle;
  \draw[->] (0,0,0)--(1.56,0,0) node[below left] {$z^{(1)}$};
  \draw[->] (0,0,0)--(0,1.44,0) node[right] {$z^{(2)}$};
  \draw[->] (0,0,0)--(0,0,1.38) node[above] {$z^{(3)}$};
  \node[blue!65!black] at (-.37,.39,.98) {预算平面 $\iota^{\mathsf T}z=1$};

  % 沿 z_0+t(z_1-z_0) 取 t=-.5 和 t=1.6 作可见端点。
  \draw[frontier] (.985,-.04,.055)--(-.212,.632,.58);
  \fill (.70,.12,.18) circle(1.7pt) node[above left] {$z_0$};
  \fill (.415,.28,.305) circle(1.7pt) node[above left] {$z_{\mathrm{MVP}}$};
  \fill (.13,.44,.43) circle(1.7pt) node[below right] {$z_1$};
  \node at (0,0,-.61) {(b) 三维权重空间中的视图};
 \end{scope}
\end{tikzpicture}
